\begin{theorem}[7.1.8]
\label{IV.2.s7.1.8-fr}
Sea $A$ un anillo local noetheriano. Las siguientes
condiciones son equivalentes:
\begin{enumerate}
\item[\textup{a)}] Todo anillo cociente íntegro de $A$ es
formalmente equidimensional.
\item[\textup{b)}] Los anillos cociente de $A$ por sus ideales
primos minimales son formalmente equidimensionales.
\item[\textup{c)}] Todo anillo local $B$ que sea una
$A$-álgebra esencialmente de tipo finito \textup{(1.3.8)}
y que sea equidimensional es formalmente
equidimensional.
\end{enumerate}
\end{theorem}

\begin{proof}
Es inmediato que \emph{c)} implica \emph{a)} y que \emph{a)} implica \emph{b)}.
Por otra parte, como todo ideal primo de $A$ contiene
un ideal primo minimal de $A$, \emph{b)} implica \emph{a)} en virtud
de (7.1.4, (ii)). Queda por probar que \emph{a)} implica \emph{c)}.

Basta demostrar que los cocientes de $B$ por sus ideales
primos minimales son formalmente equidimensionales
(7.1.2). Como todo anillo cociente de $B$ sigue siendo
una $A$-álgebra esencialmente de tipo finito (1.3.9),
se puede suponer, en primer lugar, que $B$ es
\emph{íntegro}. Si $\mathfrak q$ es la imagen inversa
en $A$ del ideal maximal de $B$, se sabe que $B$ es
también una $A_{\mathfrak q}$-álgebra esencialmente
de tipo finito (1.3.10), y de \emph{a)} y (7.1.7) se deduce
que todo anillo cociente íntegro de $A_{\mathfrak q}$
es formalmente equidimensional. Se puede, por tanto,
suponer que el homomorfismo $A\to B$ es \emph{local}.
El núcleo $\mathfrak r$ de este homomorfismo es
entonces un ideal primo de $A$; por \emph{a)}, todo anillo
cociente íntegro de $A/\mathfrak r$ es formalmente
equidimensional. Como $B$ es una
$(A/\mathfrak r)$-álgebra esencialmente de tipo
finito, se ve que también se puede suponer

\oldpage[IV-2]{186}
que $A$ es \emph{íntegro} y es un subanillo de $B$.
Se sabe (1.3.11) que $B$ es cociente de un anillo
local de la forma $C_{\mathfrak p}$, donde
$C=A[T_1,\ldots,T_n]$ es un álgebra de polinomios
y $\mathfrak p$ es un ideal primo de $C$ situado
sobre el ideal maximal $\mathfrak m$ de $A$.
En virtud de (7.1.7), basta probar que
$C_{\mathfrak p}$ es formalmente equidimensional;
es decir, es posible limitarse al caso
$B=C_{\mathfrak p}$.

Pongamos $A'=\widehat A$ y
$C'=A'[T_1,\ldots,T_n]=C\otimes_AA'$.
Hay un único ideal primo $\mathfrak p'$ de $C'$
situado sobre $\mathfrak p$, y por tanto sobre el
ideal maximal $\mathfrak mA'$ de $A'$; pongamos
$B'=C'_{\mathfrak p'}$. El homomorfismo
$B\to B'$ es local y hace de $B'$ un $B$-módulo
plano. Se sabe entonces que $\widehat{B}'$ es
un $\widehat B$-módulo plano (Bourbaki,
\emph{Alg. comm.}, cap. III, § 5,
n.\textsuperscript{o} 4, prop. 4).
Como $\widehat{B}'$ es catenario (5.6.4),
bastará probar que $\widehat{B}'$ es
\emph{equidimensional} para deducir que
$\widehat B$ también lo es (7.1.3, (i)),
lo que terminará la demostración.

Ahora bien, $A'$ es cociente de un anillo regular
por el teorema de Cohen (\textbf{0}, 19.8.8), y por tanto
también lo es $B'$ (\textbf 0, 17.3.9). En virtud de
(7.1.5), para demostrar que $B'$ es formalmente
equidimensional basta probar que $B'$ es
\emph{equidimensional}; para ello basta mostrar que
$C'$ es equidimensional, ya que $C'$ es cociente
de un anillo regular y, por tanto, catenario
(5.6.4). Los ideales primos minimales de $C'$
son los ideales $\mathfrak q'C'$, donde
$\mathfrak q'$ recorre los ideales primos minimales
de $A'$ (5.5.3), y se tiene
$C'/\mathfrak q'C'=(A'/\mathfrak q')[T_1,\ldots,T_n]$.
Como $A'$ es equidimensional por hipótesis
(puesto que $A$ es íntegro), lo mismo ocurre
con $C'$ por (5.5.4), como queríamos demostrar.
\end{proof}
